% Hieu
\section{Hiện thực cấu trúc bảng}

Phần này trình bày toàn bộ mã lệnh SQL hiện thực hóa cấu trúc vật lý của 19 bảng trong cơ sở dữ liệu \texttt{DataCenterManagement\_DB} trên hệ quản trị \textbf{Microsoft SQL Server}. Các bảng được tổ chức theo 4 phân hệ chức năng đã thiết kế tại Chương 3, bao gồm: định nghĩa kiểu dữ liệu, khóa chính (\texttt{PRIMARY KEY}), khóa ngoại (\texttt{FOREIGN KEY REFERENCES}), và ràng buộc \texttt{NOT NULL}.

\vspace{0.5em}
\noindent Trước tiên, tạo cơ sở dữ liệu và chuyển ngữ cảnh sử dụng:

\begin{tcolorbox}[colback=gray!5, colframe=gray!60, title={\textbf{Khởi tạo Database}}, fonttitle=\bfseries\small, boxrule=0.5pt, arc=2pt]
\begin{lstlisting}[language=SQL, basicstyle=\ttfamily\small, breaklines=true, showstringspaces=false, keywordstyle=\bfseries\color{blue!70!black}, stringstyle=\color{red!70!black}, commentstyle=\itshape\color{green!50!black}]
CREATE DATABASE DataCenterManagement_DB;
GO

USE DataCenterManagement_DB;
GO
\end{lstlisting}
\end{tcolorbox}

%% =============================================
%% 4.1.1 PHAN HE HA TANG & KHONG GIAN VAT LY
%% =============================================
\subsection{Phân hệ Hạ tầng \& Không gian Vật lý}

Phân hệ này gồm 5 bảng mô tả cấu trúc phân cấp vật lý của trung tâm dữ liệu: \texttt{Site} $\rightarrow$ \texttt{Room} $\rightarrow$ \texttt{Rack} $\rightarrow$ \texttt{Server} $\rightarrow$ \texttt{Workload}. Mỗi bảng con tham chiếu khóa ngoại đến bảng cha, thể hiện quan hệ chứa -- thuộc.

\vspace{0.5em}
\noindent\textbf{Bảng Site} -- Lưu trữ thông tin các trung tâm dữ liệu vật lý (Data Center Site).

\begin{tcolorbox}[colback=gray!5, colframe=gray!60, title={\textbf{CREATE TABLE Site}}, fonttitle=\bfseries\small, boxrule=0.5pt, arc=2pt]
\begin{lstlisting}[language=SQL, basicstyle=\ttfamily\small, breaklines=true, showstringspaces=false, keywordstyle=\bfseries\color{blue!70!black}, stringstyle=\color{red!70!black}, commentstyle=\itshape\color{green!50!black}]
CREATE TABLE Site (
    SiteID       VARCHAR(36) PRIMARY KEY,
    SiteName     VARCHAR(50) NOT NULL,
    Location     VARCHAR(255),
    Description  TEXT
);
\end{lstlisting}
\end{tcolorbox}

\begin{itemize}[leftmargin=2em]
    \item \texttt{SiteID}: Khóa chính, định danh duy nhất cho mỗi trung tâm dữ liệu.
    \item \texttt{SiteName}: Tên trung tâm dữ liệu, bắt buộc nhập (\texttt{NOT NULL}).
    \item \texttt{Location}: Địa chỉ vị trí địa lý, cho phép \texttt{NULL}.
    \item \texttt{Description}: Mô tả chi tiết bổ sung, kiểu \texttt{TEXT}.
\end{itemize}

\vspace{0.5em}
\noindent\textbf{Bảng Room} -- Quản lý các phòng máy chuyên dụng thuộc một Site.

\begin{tcolorbox}[colback=gray!5, colframe=gray!60, title={\textbf{CREATE TABLE Room}}, fonttitle=\bfseries\small, boxrule=0.5pt, arc=2pt]
\begin{lstlisting}[language=SQL, basicstyle=\ttfamily\small, breaklines=true, showstringspaces=false, keywordstyle=\bfseries\color{blue!70!black}, stringstyle=\color{red!70!black}, commentstyle=\itshape\color{green!50!black}]
CREATE TABLE Room (
    RoomID     VARCHAR(36) PRIMARY KEY,
    SiteID     VARCHAR(36) FOREIGN KEY REFERENCES Site(SiteID),
    RoomName   VARCHAR(30) NOT NULL,
    FloorLevel VARCHAR(10)
);
\end{lstlisting}
\end{tcolorbox}

\begin{itemize}[leftmargin=2em]
    \item \texttt{RoomID}: Khóa chính, định danh phòng máy.
    \item \texttt{SiteID}: Khóa ngoại tham chiếu đến \texttt{Site(SiteID)}, xác định phòng thuộc trung tâm dữ liệu nào.
    \item \texttt{RoomName}: Tên phòng máy, bắt buộc nhập.
    \item \texttt{FloorLevel}: Tầng đặt phòng máy.
\end{itemize}

\vspace{0.5em}
\noindent\textbf{Bảng Rack} -- Quản lý các tủ rack chứa thiết bị trong phòng máy.

\begin{tcolorbox}[colback=gray!5, colframe=gray!60, title={\textbf{CREATE TABLE Rack}}, fonttitle=\bfseries\small, boxrule=0.5pt, arc=2pt]
\begin{lstlisting}[language=SQL, basicstyle=\ttfamily\small, breaklines=true, showstringspaces=false, keywordstyle=\bfseries\color{blue!70!black}, stringstyle=\color{red!70!black}, commentstyle=\itshape\color{green!50!black}]
CREATE TABLE Rack (
    RackID         VARCHAR(36) PRIMARY KEY,
    RoomID         VARCHAR(36) FOREIGN KEY REFERENCES Room(RoomID),
    RackName       VARCHAR(30) NOT NULL,
    MaxUnit        INT,
    LocationInRoom VARCHAR(50)
);
\end{lstlisting}
\end{tcolorbox}

\begin{itemize}[leftmargin=2em]
    \item \texttt{RackID}: Khóa chính, định danh tủ rack.
    \item \texttt{RoomID}: Khóa ngoại tham chiếu đến \texttt{Room(RoomID)}.
    \item \texttt{RackName}: Tên tủ rack, bắt buộc nhập.
    \item \texttt{MaxUnit}: Số đơn vị U tối đa mà rack hỗ trợ (kiểu \texttt{INT}).
    \item \texttt{LocationInRoom}: Vị trí đặt rack trong phòng máy.
\end{itemize}

\vspace{0.5em}
\noindent\textbf{Bảng Server} -- Quản lý máy chủ vật lý lắp đặt trong tủ Rack.

\begin{tcolorbox}[colback=gray!5, colframe=gray!60, title={\textbf{CREATE TABLE Server}}, fonttitle=\bfseries\small, boxrule=0.5pt, arc=2pt]
\begin{lstlisting}[language=SQL, basicstyle=\ttfamily\small, breaklines=true, showstringspaces=false, keywordstyle=\bfseries\color{blue!70!black}, stringstyle=\color{red!70!black}, commentstyle=\itshape\color{green!50!black}]
CREATE TABLE Server (
    ServerID    VARCHAR(36) PRIMARY KEY,
    RackID      VARCHAR(36) FOREIGN KEY REFERENCES Rack(RackID),
    ServerName  VARCHAR(40) NOT NULL,
    IPAddress   VARCHAR(45),
    QRCodeStamp VARCHAR(50),
    UPosition   INT,
    Status      VARCHAR(20),
    LastSeen    DATETIME
);
\end{lstlisting}
\end{tcolorbox}

\begin{itemize}[leftmargin=2em]
    \item \texttt{ServerID}: Khóa chính, định danh máy chủ.
    \item \texttt{RackID}: Khóa ngoại tham chiếu đến \texttt{Rack(RackID)}.
    \item \texttt{ServerName}: Tên máy chủ, bắt buộc nhập.
    \item \texttt{IPAddress}: Địa chỉ IP quản lý, hỗ trợ cả IPv4/IPv6 (\texttt{VARCHAR(45)}).
    \item \texttt{QRCodeStamp}: Mã QR/ArUco Marker dán trên thiết bị vật lý, phục vụ nhận diện AR.
    \item \texttt{UPosition}: Vị trí đơn vị U trong rack (kiểu \texttt{INT}).
    \item \texttt{Status}: Trạng thái hoạt động hiện tại của máy chủ.
    \item \texttt{LastSeen}: Thời điểm hệ thống ghi nhận máy chủ trực tuyến lần cuối.
\end{itemize}

\vspace{0.5em}
\noindent\textbf{Bảng Workload} -- Quản lý các dịch vụ ảo hóa, container chạy trên máy chủ.

\begin{tcolorbox}[colback=gray!5, colframe=gray!60, title={\textbf{CREATE TABLE Workload}}, fonttitle=\bfseries\small, boxrule=0.5pt, arc=2pt]
\begin{lstlisting}[language=SQL, basicstyle=\ttfamily\small, breaklines=true, showstringspaces=false, keywordstyle=\bfseries\color{blue!70!black}, stringstyle=\color{red!70!black}, commentstyle=\itshape\color{green!50!black}]
CREATE TABLE Workload (
    WorkloadID   VARCHAR(36) PRIMARY KEY,
    ServerID     VARCHAR(36) FOREIGN KEY REFERENCES Server(ServerID),
    WorkloadName VARCHAR(50) NOT NULL,
    ContainerID  INT,
    ServiceType  VARCHAR(40),
    Status       VARCHAR(20)
);
\end{lstlisting}
\end{tcolorbox}

\begin{itemize}[leftmargin=2em]
    \item \texttt{WorkloadID}: Khóa chính, định danh khối tải công việc.
    \item \texttt{ServerID}: Khóa ngoại tham chiếu đến \texttt{Server(ServerID)}.
    \item \texttt{WorkloadName}: Tên dịch vụ/container, bắt buộc nhập.
    \item \texttt{ContainerID}: Mã định danh container nội bộ (kiểu \texttt{INT}).
    \item \texttt{ServiceType}: Loại dịch vụ (Web Proxy, Database, Microservice, Cache...).
    \item \texttt{Status}: Trạng thái hoạt động của workload.
\end{itemize}

%% =============================================
%% 4.1.2 PHAN HE GIAM SAT & VAN HANH
%% =============================================
\subsection{Phân hệ Giám sát \& Vận hành Hệ thống}

Phân hệ này gồm 6 bảng phục vụ giám sát hiệu năng, quản lý cảnh báo, xử lý sự cố và bảo trì thiết bị: \texttt{TelemetryMetric}, \texttt{AlertThreshold}, \texttt{Alert}, \texttt{Incident}, \texttt{MaintainTicket}, \texttt{MaintenanceLog}.

\vspace{0.5em}
\noindent\textbf{Bảng TelemetryMetric} -- Ghi nhận chỉ số hiệu năng hệ thống theo thời gian thực.

\begin{tcolorbox}[colback=gray!5, colframe=gray!60, title={\textbf{CREATE TABLE TelemetryMetric}}, fonttitle=\bfseries\small, boxrule=0.5pt, arc=2pt]
\begin{lstlisting}[language=SQL, basicstyle=\ttfamily\small, breaklines=true, showstringspaces=false, keywordstyle=\bfseries\color{blue!70!black}, stringstyle=\color{red!70!black}, commentstyle=\itshape\color{green!50!black}]
CREATE TABLE TelemetryMetric (
    MetricID        VARCHAR(36) PRIMARY KEY,
    ServerID        VARCHAR(36) FOREIGN KEY REFERENCES Server(ServerID),
    CPUUsage        FLOAT,
    RAMUsage        FLOAT,
    NetworkTraffic  FLOAT,
    ContainerStates VARCHAR(50)
);
\end{lstlisting}
\end{tcolorbox}

\begin{itemize}[leftmargin=2em]
    \item \texttt{MetricID}: Khóa chính, định danh bản ghi đo lường.
    \item \texttt{ServerID}: Khóa ngoại tham chiếu đến \texttt{Server(ServerID)}, xác định máy chủ được giám sát.
    \item \texttt{CPUUsage}, \texttt{RAMUsage}: Tỷ lệ sử dụng CPU và RAM (kiểu \texttt{FLOAT}, đơn vị \%).
    \item \texttt{NetworkTraffic}: Lưu lượng mạng (kiểu \texttt{FLOAT}, đơn vị Mbps).
    \item \texttt{ContainerStates}: Trạng thái tổng hợp của các container trên máy chủ.
\end{itemize}

\vspace{0.5em}
\noindent\textbf{Bảng AlertThreshold} -- Cấu hình các ngưỡng cảnh báo động cho các chỉ số hiệu năng.

\begin{tcolorbox}[colback=gray!5, colframe=gray!60, title={\textbf{CREATE TABLE AlertThreshold}}, fonttitle=\bfseries\small, boxrule=0.5pt, arc=2pt]
\begin{lstlisting}[language=SQL, basicstyle=\ttfamily\small, breaklines=true, showstringspaces=false, keywordstyle=\bfseries\color{blue!70!black}, stringstyle=\color{red!70!black}, commentstyle=\itshape\color{green!50!black}]
CREATE TABLE AlertThreshold (
    ThresholdID   VARCHAR(36) PRIMARY KEY,
    MetricType    VARCHAR(50),
    WarningValue  FLOAT,
    CriticalValue FLOAT,
    IsActive      BIT
);
\end{lstlisting}
\end{tcolorbox}

\begin{itemize}[leftmargin=2em]
    \item \texttt{ThresholdID}: Khóa chính, định danh cấu hình ngưỡng.
    \item \texttt{MetricType}: Loại chỉ số áp dụng ngưỡng (CPU, RAM, Network).
    \item \texttt{WarningValue}: Giá trị ngưỡng cảnh báo mức Warning (kiểu \texttt{FLOAT}).
    \item \texttt{CriticalValue}: Giá trị ngưỡng cảnh báo mức Critical (kiểu \texttt{FLOAT}).
    \item \texttt{IsActive}: Cờ kích hoạt ngưỡng (kiểu \texttt{BIT}: 0 = tắt, 1 = bật).
\end{itemize}

\vspace{0.5em}
\noindent\textbf{Bảng Alert} -- Ghi nhận các sự kiện cảnh báo phát sinh khi chỉ số vượt ngưỡng.

\begin{tcolorbox}[colback=gray!5, colframe=gray!60, title={\textbf{CREATE TABLE Alert}}, fonttitle=\bfseries\small, boxrule=0.5pt, arc=2pt]
\begin{lstlisting}[language=SQL, basicstyle=\ttfamily\small, breaklines=true, showstringspaces=false, keywordstyle=\bfseries\color{blue!70!black}, stringstyle=\color{red!70!black}, commentstyle=\itshape\color{green!50!black}]
CREATE TABLE Alert (
    AlertID        VARCHAR(36) PRIMARY KEY,
    ServerID       VARCHAR(36) FOREIGN KEY REFERENCES Server(ServerID),
    ThresholdID    VARCHAR(36) FOREIGN KEY REFERENCES
                       AlertThreshold(ThresholdID),
    AcknowledgedBy VARCHAR(36) FOREIGN KEY REFERENCES [User](UserID),
    Alerttype      VARCHAR(30),
    Severity       VARCHAR(15),
    State          VARCHAR(15),
    Message        TEXT,
    TriggeredAt    DATETIME,
    ResolvedAt     DATETIME
);
\end{lstlisting}
\end{tcolorbox}

\begin{itemize}[leftmargin=2em]
    \item \texttt{AlertID}: Khóa chính, định danh sự kiện cảnh báo.
    \item \texttt{ServerID}: Khóa ngoại -- máy chủ phát sinh cảnh báo.
    \item \texttt{ThresholdID}: Khóa ngoại -- cấu hình ngưỡng đã kích hoạt cảnh báo.
    \item \texttt{AcknowledgedBy}: Khóa ngoại -- người dùng xác nhận cảnh báo (cho phép \texttt{NULL} khi chưa xác nhận).
    \item \texttt{Alerttype}: Loại cảnh báo (CPU Spike, RAM High, Network Peak...).
    \item \texttt{Severity}: Mức độ nghiêm trọng (\texttt{Warning}, \texttt{Critical}).
    \item \texttt{State}: Trạng thái vòng đời cảnh báo (\texttt{Open} $\rightarrow$ \texttt{Acknowledged} $\rightarrow$ \texttt{Resolved} $\rightarrow$ \texttt{Closed}).
    \item \texttt{Message}: Nội dung chi tiết thông điệp cảnh báo.
    \item \texttt{TriggeredAt}, \texttt{ResolvedAt}: Thời điểm phát sinh và xử lý xong cảnh báo.
\end{itemize}

\vspace{0.5em}
\noindent\textbf{Bảng Incident} -- Quản lý các sự cố nghiêm trọng cần xử lý cấp thiết.

\begin{tcolorbox}[colback=gray!5, colframe=gray!60, title={\textbf{CREATE TABLE Incident}}, fonttitle=\bfseries\small, boxrule=0.5pt, arc=2pt]
\begin{lstlisting}[language=SQL, basicstyle=\ttfamily\small, breaklines=true, showstringspaces=false, keywordstyle=\bfseries\color{blue!70!black}, stringstyle=\color{red!70!black}, commentstyle=\itshape\color{green!50!black}]
CREATE TABLE Incident (
    IncidentID  VARCHAR(36) PRIMARY KEY,
    CreatedID   VARCHAR(36) FOREIGN KEY REFERENCES [User](UserID),
    Title       VARCHAR(100) NOT NULL,
    Description TEXT,
    Severity    VARCHAR(20),
    Status      VARCHAR(20),
    CreatedAt   DATETIME,
    ResolvedAt  DATETIME
);
\end{lstlisting}
\end{tcolorbox}

\begin{itemize}[leftmargin=2em]
    \item \texttt{IncidentID}: Khóa chính, định danh sự cố.
    \item \texttt{CreatedID}: Khóa ngoại tham chiếu đến \texttt{User(UserID)} -- người tạo báo cáo sự cố.
    \item \texttt{Title}: Tiêu đề mô tả ngắn gọn sự cố, bắt buộc nhập.
    \item \texttt{Description}: Mô tả chi tiết sự cố (kiểu \texttt{TEXT}).
    \item \texttt{Severity}: Mức độ nghiêm trọng (\texttt{Low}, \texttt{Medium}, \texttt{High}, \texttt{Critical}).
    \item \texttt{Status}: Trạng thái xử lý sự cố (\texttt{Open}, \texttt{In Progress}, \texttt{Resolved}, \texttt{Closed}).
    \item \texttt{CreatedAt}, \texttt{ResolvedAt}: Thời điểm tạo và xử lý xong sự cố.
\end{itemize}

\vspace{0.5em}
\noindent\textbf{Bảng MaintainTicket} -- Lập phiếu và theo dõi vòng đời bảo trì thiết bị.

\begin{tcolorbox}[colback=gray!5, colframe=gray!60, title={\textbf{CREATE TABLE MaintainTicket}}, fonttitle=\bfseries\small, boxrule=0.5pt, arc=2pt]
\begin{lstlisting}[language=SQL, basicstyle=\ttfamily\small, breaklines=true, showstringspaces=false, keywordstyle=\bfseries\color{blue!70!black}, stringstyle=\color{red!70!black}, commentstyle=\itshape\color{green!50!black}]
CREATE TABLE MaintainTicket (
    TicketID    VARCHAR(36) PRIMARY KEY,
    IncidentID  VARCHAR(36) FOREIGN KEY REFERENCES
                    Incident(IncidentID),
    ServerID    VARCHAR(36) FOREIGN KEY REFERENCES
                    Server(ServerID),
    AssignedTo  VARCHAR(36) FOREIGN KEY REFERENCES
                    [User](UserID),
    CreatedBy   VARCHAR(36) FOREIGN KEY REFERENCES
                    [User](UserID),
    Priority    VARCHAR(20),
    Status      VARCHAR(20),
    ScheduledAt DATETIME,
    StartedAt   DATETIME,
    CompletedAt DATETIME
);
\end{lstlisting}
\end{tcolorbox}

\begin{itemize}[leftmargin=2em]
    \item \texttt{TicketID}: Khóa chính, định danh phiếu bảo trì.
    \item \texttt{IncidentID}: Khóa ngoại -- sự cố liên quan (cho phép \texttt{NULL} nếu bảo trì định kỳ).
    \item \texttt{ServerID}: Khóa ngoại -- máy chủ cần bảo trì.
    \item \texttt{AssignedTo}: Khóa ngoại -- kỹ thuật viên được phân công xử lý.
    \item \texttt{CreatedBy}: Khóa ngoại -- người tạo phiếu bảo trì.
    \item \texttt{Priority}: Mức độ ưu tiên (\texttt{Low}, \texttt{Medium}, \texttt{High}, \texttt{Critical}).
    \item \texttt{Status}: Trạng thái phiếu bảo trì (\texttt{Open}, \texttt{In Progress}, \texttt{Completed}).
    \item \texttt{ScheduledAt}, \texttt{StartedAt}, \texttt{CompletedAt}: Chuỗi thời gian theo dõi vòng đời xử lý phiếu.
\end{itemize}

\vspace{0.5em}
\noindent\textbf{Bảng MaintenanceLog} -- Chi tiết nhật ký xử lý kỹ thuật và thao tác hỗ trợ AR.

\begin{tcolorbox}[colback=gray!5, colframe=gray!60, title={\textbf{CREATE TABLE MaintenanceLog}}, fonttitle=\bfseries\small, boxrule=0.5pt, arc=2pt]
\begin{lstlisting}[language=SQL, basicstyle=\ttfamily\small, breaklines=true, showstringspaces=false, keywordstyle=\bfseries\color{blue!70!black}, stringstyle=\color{red!70!black}, commentstyle=\itshape\color{green!50!black}]
CREATE TABLE MaintenanceLog (
    LogID              VARCHAR(36) PRIMARY KEY,
    TicketID           VARCHAR(36) FOREIGN KEY REFERENCES
                           MaintainTicket(TicketID),
    TechnicianID       VARCHAR(36) FOREIGN KEY REFERENCES
                           [User](UserID),
    AssetID            VARCHAR(36) FOREIGN KEY REFERENCES
                           Asset(AssetID),
    ARSessionUsed      VARCHAR(50),
    InvestigationNotes TEXT,
    RootCause          TEXT,
    CorrectiveAction   TEXT,
    LoggedAt           DATETIME
);
\end{lstlisting}
\end{tcolorbox}

\begin{itemize}[leftmargin=2em]
    \item \texttt{LogID}: Khóa chính, định danh nhật ký xử lý.
    \item \texttt{TicketID}: Khóa ngoại -- phiếu bảo trì liên quan.
    \item \texttt{TechnicianID}: Khóa ngoại -- kỹ thuật viên thực hiện.
    \item \texttt{AssetID}: Khóa ngoại -- tài sản được xử lý.
    \item \texttt{ARSessionUsed}: Mã phiên sử dụng kính AR trong quá trình xử lý (nếu có).
    \item \texttt{InvestigationNotes}: Ghi chú điều tra ban đầu.
    \item \texttt{RootCause}: Nguyên nhân gốc rễ xác định được.
    \item \texttt{CorrectiveAction}: Hành động khắc phục đã thực hiện.
    \item \texttt{LoggedAt}: Thời điểm ghi nhật ký.
\end{itemize}

%% =============================================
%% 4.1.3 PHAN HE QUAN LY TAI SAN & BAO HANH
%% =============================================
\subsection{Phân hệ Quản lý Tài sản \& Bảo hành}

Phân hệ này gồm 4 bảng quản lý vòng đời tài sản phần cứng: \texttt{AssetCategory}, \texttt{Asset}, \texttt{Warranty}, và \texttt{AssetStatusChange}.

\vspace{0.5em}
\noindent\textbf{Bảng AssetCategory} -- Danh mục phân loại tài sản phần cứng.

\begin{tcolorbox}[colback=gray!5, colframe=gray!60, title={\textbf{CREATE TABLE AssetCategory}}, fonttitle=\bfseries\small, boxrule=0.5pt, arc=2pt]
\begin{lstlisting}[language=SQL, basicstyle=\ttfamily\small, breaklines=true, showstringspaces=false, keywordstyle=\bfseries\color{blue!70!black}, stringstyle=\color{red!70!black}, commentstyle=\itshape\color{green!50!black}]
CREATE TABLE AssetCategory (
    CategoryID   VARCHAR(36) PRIMARY KEY,
    CategoryName VARCHAR(50) NOT NULL,
    CategoryCode VARCHAR(20),
    Description  TEXT
);
\end{lstlisting}
\end{tcolorbox}

\begin{itemize}[leftmargin=2em]
    \item \texttt{CategoryID}: Khóa chính, định danh danh mục tài sản.
    \item \texttt{CategoryName}: Tên danh mục (Rack Server, Network Switch...), bắt buộc nhập.
    \item \texttt{CategoryCode}: Mã viết tắt danh mục (VD: \texttt{SRV-RACK}, \texttt{NET-SW}).
    \item \texttt{Description}: Mô tả chi tiết danh mục.
\end{itemize}

\vspace{0.5em}
\noindent\textbf{Bảng Asset} -- Quản lý thông tin chi tiết từng tài sản phần cứng.

\begin{tcolorbox}[colback=gray!5, colframe=gray!60, title={\textbf{CREATE TABLE Asset}}, fonttitle=\bfseries\small, boxrule=0.5pt, arc=2pt]
\begin{lstlisting}[language=SQL, basicstyle=\ttfamily\small, breaklines=true, showstringspaces=false, keywordstyle=\bfseries\color{blue!70!black}, stringstyle=\color{red!70!black}, commentstyle=\itshape\color{green!50!black}]
CREATE TABLE Asset (
    AssetID         VARCHAR(36) PRIMARY KEY,
    CategoryID      VARCHAR(36) FOREIGN KEY REFERENCES
                        AssetCategory(CategoryID),
    ServerID        VARCHAR(36) FOREIGN KEY REFERENCES
                        Server(ServerID),
    AssetTag        VARCHAR(20),
    AssetName       VARCHAR(50) NOT NULL,
    SerialNumber    VARCHAR(50),
    Manufacturer    VARCHAR(40),
    Model           VARCHAR(40),
    CPUModel        VARCHAR(40),
    RAMTotalGB      INT,
    AcquisitionDate DATETIME,
    AcquisitionCost FLOAT,
    Status          VARCHAR(20)
);
\end{lstlisting}
\end{tcolorbox}

\begin{itemize}[leftmargin=2em]
    \item \texttt{AssetID}: Khóa chính, định danh tài sản phần cứng.
    \item \texttt{CategoryID}: Khóa ngoại -- phân loại tài sản thuộc danh mục nào.
    \item \texttt{ServerID}: Khóa ngoại -- máy chủ mà tài sản đang gắn kết (cho phép \texttt{NULL} nếu tài sản ở kho).
    \item \texttt{AssetTag}: Mã thẻ quản lý tài sản nội bộ.
    \item \texttt{AssetName}: Tên tài sản, bắt buộc nhập (VD: Dell PowerEdge R750).
    \item \texttt{SerialNumber}: Số sê-ri phần cứng từ nhà sản xuất.
    \item \texttt{Manufacturer}: Hãng sản xuất (Dell, HPE, Cisco...).
    \item \texttt{Model}: Model cụ thể của thiết bị.
    \item \texttt{CPUModel}: Model CPU trang bị (nếu là máy chủ).
    \item \texttt{RAMTotalGB}: Tổng dung lượng RAM (kiểu \texttt{INT}, đơn vị GB).
    \item \texttt{AcquisitionDate}: Ngày mua/tiếp nhận tài sản.
    \item \texttt{AcquisitionCost}: Giá trị mua (kiểu \texttt{FLOAT}, đơn vị USD).
    \item \texttt{Status}: Trạng thái tài sản hiện tại (In Use, In Storage...).
\end{itemize}

\vspace{0.5em}
\noindent\textbf{Bảng Warranty} -- Theo dõi thời hạn bảo hành tài sản với nhà cung cấp.

\begin{tcolorbox}[colback=gray!5, colframe=gray!60, title={\textbf{CREATE TABLE Warranty}}, fonttitle=\bfseries\small, boxrule=0.5pt, arc=2pt]
\begin{lstlisting}[language=SQL, basicstyle=\ttfamily\small, breaklines=true, showstringspaces=false, keywordstyle=\bfseries\color{blue!70!black}, stringstyle=\color{red!70!black}, commentstyle=\itshape\color{green!50!black}]
CREATE TABLE Warranty (
    WarrantyID     VARCHAR(36) PRIMARY KEY,
    AssetID        VARCHAR(36) FOREIGN KEY REFERENCES
                       Asset(AssetID),
    WarrantyType   VARCHAR(20),
    WarrantyName   VARCHAR(50),
    ContractNumber VARCHAR(30),
    StartDate      DATETIME,
    EndDate        DATETIME,
    AlertBeforeDays INT,
    Status         VARCHAR(20)
);
\end{lstlisting}
\end{tcolorbox}

\begin{itemize}[leftmargin=2em]
    \item \texttt{WarrantyID}: Khóa chính, định danh hợp đồng bảo hành.
    \item \texttt{AssetID}: Khóa ngoại -- tài sản được bảo hành.
    \item \texttt{WarrantyType}: Loại bảo hành (Vendor, Extended...).
    \item \texttt{WarrantyName}: Tên gói bảo hành (VD: Dell ProSupport 3 Years).
    \item \texttt{ContractNumber}: Số hợp đồng bảo hành.
    \item \texttt{StartDate}, \texttt{EndDate}: Thời gian hiệu lực bảo hành.
    \item \texttt{AlertBeforeDays}: Số ngày cảnh báo trước khi hết hạn bảo hành (kiểu \texttt{INT}).
    \item \texttt{Status}: Trạng thái bảo hành (\texttt{Active}, \texttt{Expired}, \texttt{Claimed}).
\end{itemize}

\vspace{0.5em}
\noindent\textbf{Bảng AssetStatusChange} -- Lịch sử biến động trạng thái tài sản theo thời gian.

\begin{tcolorbox}[colback=gray!5, colframe=gray!60, title={\textbf{CREATE TABLE AssetStatusChange}}, fonttitle=\bfseries\small, boxrule=0.5pt, arc=2pt]
\begin{lstlisting}[language=SQL, basicstyle=\ttfamily\small, breaklines=true, showstringspaces=false, keywordstyle=\bfseries\color{blue!70!black}, stringstyle=\color{red!70!black}, commentstyle=\itshape\color{green!50!black}]
CREATE TABLE AssetStatusChange (
    ChangeID       VARCHAR(36) PRIMARY KEY,
    AssetID        VARCHAR(36) FOREIGN KEY REFERENCES
                       Asset(AssetID),
    PreviousStatus VARCHAR(15),
    NewStatus      VARCHAR(15),
    Reason         TEXT,
    ChangedAt      DATETIME
);
\end{lstlisting}
\end{tcolorbox}

\begin{itemize}[leftmargin=2em]
    \item \texttt{ChangeID}: Khóa chính, định danh bản ghi thay đổi trạng thái.
    \item \texttt{AssetID}: Khóa ngoại -- tài sản có trạng thái thay đổi.
    \item \texttt{PreviousStatus}: Trạng thái trước khi thay đổi.
    \item \texttt{NewStatus}: Trạng thái sau khi thay đổi.
    \item \texttt{Reason}: Lý do thay đổi trạng thái (kiểu \texttt{TEXT}).
    \item \texttt{ChangedAt}: Thời điểm ghi nhận thay đổi.
\end{itemize}

%% =============================================
%% 4.1.4 PHAN HE BAO MAT, PHAN QUYEN & KIEM TOAN
%% =============================================
\subsection{Phân hệ Bảo mật, Phân quyền \& Kiểm toán}

Phân hệ này gồm 4 bảng quản lý người dùng, vai trò, phân quyền và nhật ký kiểm toán: \texttt{User}, \texttt{Role}, \texttt{UserRole}, và \texttt{AuditLog}.

\vspace{0.5em}
\noindent\textbf{Bảng User} -- Quản lý thông tin tài khoản người dùng hệ thống.

\begin{tcolorbox}[colback=gray!5, colframe=gray!60, title={\textbf{CREATE TABLE [User]}}, fonttitle=\bfseries\small, boxrule=0.5pt, arc=2pt]
\begin{lstlisting}[language=SQL, basicstyle=\ttfamily\small, breaklines=true, showstringspaces=false, keywordstyle=\bfseries\color{blue!70!black}, stringstyle=\color{red!70!black}, commentstyle=\itshape\color{green!50!black}]
CREATE TABLE [User] (
    UserID       VARCHAR(36) PRIMARY KEY,
    Username     VARCHAR(30) NOT NULL,
    PasswordHash VARCHAR(255) NOT NULL,
    FullName     VARCHAR(50) NOT NULL,
    Email        VARCHAR(50) NULL,
    Status       VARCHAR(15),
    CreatedAt    DATETIME
);
\end{lstlisting}
\end{tcolorbox}

\begin{itemize}[leftmargin=2em]
    \item \texttt{UserID}: Khóa chính, định danh tài khoản người dùng.
    \item \texttt{Username}: Tên đăng nhập, bắt buộc nhập và duy nhất trên hệ thống.
    \item \texttt{PasswordHash}: Mật khẩu đã mã hóa băm (\texttt{NOT NULL}), lưu trữ an toàn (không lưu plaintext).
    \item \texttt{FullName}: Họ và tên đầy đủ, bắt buộc nhập.
    \item \texttt{Email}: Địa chỉ email, cho phép \texttt{NULL}.
    \item \texttt{Status}: Trạng thái tài khoản (\texttt{Active}, \texttt{Inactive}, \texttt{Locked}).
    \item \texttt{CreatedAt}: Thời điểm tạo tài khoản.
\end{itemize}

\noindent\textit{Lưu ý:} Tên bảng \texttt{User} được đặt trong dấu ngoặc vuông \texttt{[User]} do \texttt{User} là từ khóa dành riêng (reserved keyword) trong SQL Server.

\vspace{0.5em}
\noindent\textbf{Bảng Role} -- Định nghĩa các vai trò chức năng trong hệ thống RBAC.

\begin{tcolorbox}[colback=gray!5, colframe=gray!60, title={\textbf{CREATE TABLE [Role]}}, fonttitle=\bfseries\small, boxrule=0.5pt, arc=2pt]
\begin{lstlisting}[language=SQL, basicstyle=\ttfamily\small, breaklines=true, showstringspaces=false, keywordstyle=\bfseries\color{blue!70!black}, stringstyle=\color{red!70!black}, commentstyle=\itshape\color{green!50!black}]
CREATE TABLE [Role] (
    RoleID      VARCHAR(36) PRIMARY KEY,
    RoleName    VARCHAR(50) NOT NULL,
    Description TEXT
);
\end{lstlisting}
\end{tcolorbox}

\begin{itemize}[leftmargin=2em]
    \item \texttt{RoleID}: Khóa chính, định danh vai trò.
    \item \texttt{RoleName}: Tên vai trò (Administrator, Technician, Viewer), bắt buộc nhập.
    \item \texttt{Description}: Mô tả quyền hạn và phạm vi trách nhiệm của vai trò.
\end{itemize}

\vspace{0.5em}
\noindent\textbf{Bảng UserRole} -- Bảng trung gian phân quyền quan hệ nhiều--nhiều ($N$--$M$) giữa User và Role.

\begin{tcolorbox}[colback=gray!5, colframe=gray!60, title={\textbf{CREATE TABLE UserRole}}, fonttitle=\bfseries\small, boxrule=0.5pt, arc=2pt]
\begin{lstlisting}[language=SQL, basicstyle=\ttfamily\small, breaklines=true, showstringspaces=false, keywordstyle=\bfseries\color{blue!70!black}, stringstyle=\color{red!70!black}, commentstyle=\itshape\color{green!50!black}]
CREATE TABLE UserRole (
    UserRoleID VARCHAR(36) PRIMARY KEY,
    UserID     VARCHAR(36) FOREIGN KEY REFERENCES [User](UserID),
    RoleID     VARCHAR(36) FOREIGN KEY REFERENCES [Role](RoleID)
);
\end{lstlisting}
\end{tcolorbox}

\begin{itemize}[leftmargin=2em]
    \item \texttt{UserRoleID}: Khóa chính, định danh bản ghi phân quyền.
    \item \texttt{UserID}: Khóa ngoại tham chiếu đến \texttt{User(UserID)}.
    \item \texttt{RoleID}: Khóa ngoại tham chiếu đến \texttt{Role(RoleID)}.
\end{itemize}

\noindent\textit{Lưu ý:} Cặp (\texttt{UserID}, \texttt{RoleID}) cần được ràng buộc \texttt{UNIQUE} để đảm bảo một người dùng không được gán cùng một vai trò hai lần.

\vspace{0.5em}
\noindent\textbf{Bảng AuditLog} -- Lưu nhật ký kiểm toán và truy vết mọi thao tác trên hệ thống.

\begin{tcolorbox}[colback=gray!5, colframe=gray!60, title={\textbf{CREATE TABLE AuditLog}}, fonttitle=\bfseries\small, boxrule=0.5pt, arc=2pt]
\begin{lstlisting}[language=SQL, basicstyle=\ttfamily\small, breaklines=true, showstringspaces=false, keywordstyle=\bfseries\color{blue!70!black}, stringstyle=\color{red!70!black}, commentstyle=\itshape\color{green!50!black}]
CREATE TABLE AuditLog (
    LogID      VARCHAR(36) PRIMARY KEY,
    IncidentID VARCHAR(36) FOREIGN KEY REFERENCES
                   Incident(IncidentID),
    UserID     VARCHAR(36) FOREIGN KEY REFERENCES
                   [User](UserID),
    Action     VARCHAR(30),
    EntityType VARCHAR(30),
    EntityID   VARCHAR(36),
    Timestamp  DATETIME,
    IPAddress  VARCHAR(15),
    Details    TEXT
);
\end{lstlisting}
\end{tcolorbox}

\begin{itemize}[leftmargin=2em]
    \item \texttt{LogID}: Khóa chính, định danh nhật ký kiểm toán.
    \item \texttt{IncidentID}: Khóa ngoại -- sự cố liên quan (cho phép \texttt{NULL} với thao tác không gắn sự cố).
    \item \texttt{UserID}: Khóa ngoại -- người dùng thực hiện thao tác.
    \item \texttt{Action}: Loại thao tác (\texttt{INSERT}, \texttt{UPDATE}, \texttt{DELETE}, \texttt{LOGIN}).
    \item \texttt{EntityType}: Loại thực thể bị tác động (Incident, User, Server...).
    \item \texttt{EntityID}: ID của thực thể bị tác động.
    \item \texttt{Timestamp}: Thời điểm thực hiện thao tác.
    \item \texttt{IPAddress}: Địa chỉ IP của người thực hiện.
    \item \texttt{Details}: Chi tiết bổ sung về thao tác (kiểu \texttt{TEXT}).
\end{itemize}

\noindent\textit{Lưu ý:} Theo quy tắc nghiệp vụ BR-10, bảng \texttt{AuditLog} chỉ cho phép thao tác \texttt{INSERT}. Mọi thao tác \texttt{UPDATE} hoặc \texttt{DELETE} sẽ bị từ chối bởi Trigger bảo vệ (trình bày tại Mục 4.4).

%% =============================================
%% 4.1.5 TONG HOP CAU TRUC BANG
%% =============================================
\subsection{Tổng hợp cấu trúc bảng toàn hệ thống}

Bảng tổng hợp dưới đây liệt kê toàn bộ 19 bảng đã hiện thực, bao gồm số lượng cột, khóa chính, khóa ngoại và hệ quản trị sử dụng.

\begin{longtable}{|c|l|c|c|c|}
\caption{Tổng hợp cấu trúc 19 bảng đã hiện thực trên SQL Server} \label{tab:tonghop_bang} \\
\hline
\textbf{STT} & \textbf{Tên bảng} & \textbf{Số cột} & \textbf{Số FK} & \textbf{Phân hệ} \\ \hline
\endfirsthead
\hline
\textbf{STT} & \textbf{Tên bảng} & \textbf{Số cột} & \textbf{Số FK} & \textbf{Phân hệ} \\ \hline
\endhead
\hline
\endfoot
1 & \texttt{Site} & 4 & 0 & Hạ tầng \\ \hline
2 & \texttt{Room} & 4 & 1 & Hạ tầng \\ \hline
3 & \texttt{Rack} & 5 & 1 & Hạ tầng \\ \hline
4 & \texttt{Server} & 8 & 1 & Hạ tầng \\ \hline
5 & \texttt{Workload} & 6 & 1 & Hạ tầng \\ \hline
6 & \texttt{TelemetryMetric} & 6 & 1 & Giám sát \\ \hline
7 & \texttt{AlertThreshold} & 5 & 0 & Giám sát \\ \hline
8 & \texttt{Alert} & 10 & 3 & Giám sát \\ \hline
9 & \texttt{Incident} & 8 & 1 & Giám sát \\ \hline
10 & \texttt{MaintainTicket} & 10 & 4 & Giám sát \\ \hline
11 & \texttt{MaintenanceLog} & 9 & 3 & Giám sát \\ \hline
12 & \texttt{AssetCategory} & 4 & 0 & Tài sản \\ \hline
13 & \texttt{Asset} & 13 & 2 & Tài sản \\ \hline
14 & \texttt{Warranty} & 9 & 1 & Tài sản \\ \hline
15 & \texttt{AssetStatusChange} & 6 & 1 & Tài sản \\ \hline
16 & \texttt{User} & 7 & 0 & Bảo mật \\ \hline
17 & \texttt{Role} & 3 & 0 & Bảo mật \\ \hline
18 & \texttt{UserRole} & 3 & 2 & Bảo mật \\ \hline
19 & \texttt{AuditLog} & 9 & 2 & Bảo mật \\ \hline
\end{longtable}

\noindent Toàn bộ 19 bảng đã được hiện thực thành công trên nền tảng \textbf{Microsoft SQL Server}, sử dụng cơ sở dữ liệu \texttt{DataCenterManagement\_DB}. Cấu trúc bảng phản ánh đúng thiết kế logic tại Chương 3 với đầy đủ khóa chính, khóa ngoại, kiểu dữ liệu và ràng buộc \texttt{NOT NULL}. Các ràng buộc bổ sung (\texttt{CHECK}, \texttt{UNIQUE}, Trigger) sẽ được trình bày chi tiết tại các mục tiếp theo của Chương 4.
